\begin{pastexam}[label={px:nchu-115-is-24}]{中興資工所乙組 115 學年度 第 24 題}
（單選；每題 2 分，答錯不倒扣）

What is the time complexity of the following code?
\begin{minted}[fontsize=\small]{c}
for (int i = 1; i < n; i *= 2) {
    for (int j = 0; j < i; j++) {   // constant work
    }
}
\end{minted}
A.~O(n) \qquad B.~O(n log n) \qquad C.~O(log n) \qquad D.~O(n$^2$)
\end{pastexam}

\begin{pastsoln}
答案是 A。

外層的 $i$ 依序是 $1, 2, 4, \ldots, 2^k$，其中 $2^k$ 是小於 $n$ 的最大的 $2$ 的次方（$n \ge 2$）。
內層在外層是 $i$ 時跑 $i$ 次，所以內層總次數是等比級數（資結書第~1.7~節）：
\[
1 + 2 + 4 + \cdots + 2^k = 2^{k+1} - 1 < 2 \cdot n.
\]
另一方面 $2^{k+1} \ge n$，總次數 $\ge n - 1$。 外層本身只跑 $k + 1 = O(\log n)$ 次。 所以總時間
$\Theta(n)$，選 A。 例如 $n = 100$：$i = 1, 2, \ldots, 64$，內層共 $127$ 次。
\begin{itemize}
\item B 錯：把外層 $\log n$ 次乘上內層「最多 $n$ 次」得到的上界 $O(n \log n)$ 不緊。 內層次數
是一直加倍的，總和由最後一項主導。
\item C 錯：那只是外層的次數。
\item D 錯：差得更遠。
\end{itemize}
嚴格說 $O(n \log n)$、$O(n^2)$ 也都是這段程式的上界（大 $O$ 只是上界），但題目問的是它的
時間複雜度，要給最緊的那一個，也就是 $\Theta(n)$ 對應的 A。
\end{pastsoln}
